\begin{introduction}

\newcommand{\introheading}[1]{%
  \par\medskip
  \noindent\textbf{#1}\par\smallskip
}

\introheading{Motivations}
La boiterie en élevage laitier constitue un enjeu majeur de bien-être animal et de
performance technico-économique. À l'échelle mondiale, elle est considérée comme l'un des
trois problèmes de santé les plus fréquents chez les vaches laitières, avec des prévalences
estimées entre 20\,\% et 55\,\% selon les études et les régions
\citep{whay2003assessment,flower2006effect}. Le coût économique est substantiel~: les
estimations publiées varient de 75 à plus de 300~USD par cas clinique, incluant des pertes
de production laitière qui peuvent atteindre 270 à 574~kg par lactation, une dégradation
de la fertilité, des augmentations de réformes anticipées et des coûts vétérinaires
directs \citep{green2002financial,bruijnis2010assessing}. Pour un troupeau de 100~vaches
soumis à une prévalence de 25\,\%, la perte annuelle attribuable aux seuls cas incidents
est estimée entre 1\,900 et 7\,500~USD. Au-delà de l'impact financier, la boiterie soulève
un enjeu éthique de bien-être animal, la douleur associée étant reconnue comme l'une des
principales sources de souffrance chronique chez les vaches laitières
\citep{whay2003assessment}.

Face à ces enjeux, l'élevage de précision (\textit{Precision Livestock Farming}, PLF)
propose d'automatiser le suivi comportemental des animaux à l'aide de capteurs portés
et d'outils numériques \citep{palma2025scoping,michelena2025plf_review}. Les
podomètres et accéléromètres aujourd'hui répandus dans les élevages laitiers génèrent
des séries temporelles continues à faible coût marginal, ouvrant la voie à des systèmes
de détection précoce intégrés au \textit{workflow} de ferme. Ce mémoire se positionne
dans cette dynamique en proposant un système de détection de boiterie probable conçu
pour être à la fois interprétable, reproductible et exploitable par l'éleveur.

\introheading{Contexte et problématique}
En pratique, l'identification précoce de la boiterie repose souvent sur des observations
humaines, utiles mais sensibles à la charge de travail et à l'expérience des observateurs.
La concordance inter-observateurs sur les grilles de \textit{scoring} locomoteur varie
significativement d'une étude à l'autre
\citep{sprecher1997lameness,flower2006effect}, ce qui motive le recours à des capteurs
comportementaux pour une détection plus continue et plus objective.

L'automatisation soulève cependant deux difficultés majeures. D'une part, les signaux
des capteurs (activité, transitions posturales, dynamique de mouvement) sont bruyants et
hétérogènes entre individus. D'autre part, un score d'anomalie n'est pas, à lui seul, une
preuve clinique suffisante de boiterie~; et l'absence fréquente de labels cliniques
synchronisés interdit de formuler le problème comme une tâche de classification supervisée
standard \citep{alsaaod2019automatic,oleary2020accelerometers}.

La littérature reflète d'ailleurs cette difficulté~: les performances rapportées varient
de 22\,\% à 91\,\% de sensibilité selon les capteurs utilisés, les définitions cliniques
retenues et les protocoles de validation
\citep{oleary2020accelerometers,gonzalezgarcia2023systematic}, ce qui complique la
comparaison directe entre travaux et la transposition d'une étude à l'autre. Plusieurs
solutions commerciales déployées en élevage (CowManager, Heatime, IceQube) fonctionnent
par ailleurs comme des boîtes noires propriétaires, rarement évaluées de manière
indépendante \citep{rutten2013invited}, ce qui motive le développement d'approches
ouvertes, documentées et auditables.

Ensemble, ces constats conduisent au problème de recherche suivant~: \textit{comment construire un
système de détection de boiterie probable, interprétable et robuste, à partir de données
capteurs comportementaux, sans prétendre au diagnostic clinique~?}

La réponse proposée repose sur une architecture hybride combinant un détecteur d'anomalies
non supervisé (Isolation Forest) et des règles métier explicites, validée à
travers un protocole en deux niveaux.

\introheading{Objectif général et objectifs spécifiques}
L'objectif général consiste à développer et valider une chaîne de traitement reproductible
permettant de détecter des épisodes de boiterie probable à partir de données capteurs
bovines, en combinant détection d'anomalies et règles métier interprétables. Cet objectif
peut se décliner en cinq sous-objectifs spécifiques~:

\begin{enumerate}
  \item Concevoir un \textit{pipeline} complet couvrant l'ensemble de la chaîne~:
  ingestion, prétraitement, extraction de variables dérivées, détection d'anomalies et
  filtrage par règles métier.
  \item Assurer l'interprétabilité opérationnelle des alertes via des règles explicites
  comme la persistance temporelle, la cohérence comportementale multi-signaux et
  l'espacement des notifications.
  \item Évaluer la robustesse du système selon deux protocoles complémentaires~:
  validation manuelle par injections contrôlées et validation robuste
  multi-configurations.
  \item Comparer le détecteur retenu (Isolation Forest) à une alternative non
  supervisée (Local Outlier Factor) dans une étude d'ablation élargie.
  \item Garantir la traçabilité et la reproductibilité des résultats par tests automatisés,
  vérification d'intégrité et gel explicite des seuils.
\end{enumerate}

\introheading{Positionnement et portée}
Le système proposé produit des \textit{alertes de boiterie probable}. Il s'agit d'une aide
à la décision pour le suivi d'élevage, et non d'un dispositif de diagnostic vétérinaire.
Un tel positionnement est volontaire et cohérent avec les limites des données disponibles
dans le cadre du mémoire, car le corpus ne dispose d'aucun étalon clinique
(\textit{gold standard}) vétérinaire synchronisé avec les enregistrements capteurs sur
l'ensemble des animaux et des périodes d'observation.

La prudence adoptée ici est conforme aux recommandations de la littérature, qui souligne la nécessité
de distinguer clairement entre systèmes de détection automatisée et outils diagnostiques
validés cliniquement \citep{alsaaod2019automatic}. Le système vise à signaler des déviations
comportementales compatibles avec une boiterie, afin de prioriser une vérification terrain.

\introheading{Méthodologie}
La démarche méthodologique articule trois volets complémentaires. Premièrement,
un pipeline de traitement modulaire en quatre étapes (chargement et filtrage des
données, extraction d'environ 50~variables dérivées, détection d'anomalies par
Isolation Forest, application de règles métier) transforme les signaux capteurs
bruts en épisodes de boiterie probable assortis d'un score de risque continu.
Deuxièmement, deux protocoles de validation complémentaires évaluent la robustesse du
système~: une validation manuelle par injections contrôlées (cas détectables, \textit{borderline}
et non-détectables) confirme la capacité du pipeline à récupérer les épisodes simulés,
et une validation robuste multi-configurations explore systématiquement l'espace des
hyperparamètres pour vérifier la stabilité de la calibration retenue. Troisièmement,
une étude d'ablation compare quatre variantes du pipeline pour isoler la contribution
de chaque composante, complétée par des tests statistiques appariés (McNemar exact et
\textit{Wilcoxon signed-rank}), des intervalles de confiance à 95\,\% par
rééchantillonnage \textit{bootstrap}, des courbes de sensibilité/spécificité et
précision/rappel, ainsi qu'une calibration isotonique des scores. L'ensemble de la
démarche s'appuie sur une traçabilité par empreintes cryptographiques SHA-256 des
artefacts expérimentaux, qui garantit la reproductibilité bit-à-bit des résultats.

\introheading{Contributions du mémoire}
Les contributions principales sont les suivantes~:
\begin{enumerate}
  \item une architecture logicielle modulaire et reproductible combinant
  Isolation Forest et règles métier explicites, couvrant toute la chaîne de
  détection depuis les données capteurs brutes jusqu'aux notifications interprétables~;
  \item un protocole de validation en deux niveaux (manuelle et robuste) couplé à une
  étude d'ablation comparant quatre variantes du pipeline (règles seules,
  Isolation Forest seul, Isolation Forest + règles, \textit{Local
  Outlier Factor} + règles), complétée par des tests statistiques appariés (McNemar
  exact, \textit{Wilcoxon signed-rank}) pour établir la significativité des écarts~;
  \item une comparaison à une \textit{baseline} pédométrique historique
  (Alsaaod 2012) et une calibration isotonique des scores, assorties d'intervalles de
  confiance à 95\,\% par bootstrap~;
  \item une couche d'interprétabilité associant analyse SHAP, courbes
  sensibilité/spécificité et précision/rappel, ainsi qu'un \textit{dashboard} Streamlit
  interactif rendant les alertes directement exploitables par l'utilisateur final~;
  \item une matrice RACI pour l'intégration opérationnelle du système dans le
  workflow de ferme, complétée par des mécanismes de traçabilité par
  empreintes cryptographiques SHA-256 garantissant l'auditabilité des résultats.
\end{enumerate}

\introheading{Structure du mémoire}
Le mémoire s'articule autour de cinq chapitres complétés par une conclusion.
Le Chapitre~\ref{chapitre:contexte} présente le contexte scientifique et la revue de
littérature sur la détection de boiterie bovine. Le Chapitre~\ref{chapitre:methodologie}
détaille les données, la méthodologie et les métriques d'évaluation. Le
Chapitre~\ref{chapitre:architecture} décrit l'architecture logicielle et les choix
d'implémentation, incluant les détails algorithmiques d'Isolation Forest et de
Local Outlier Factor. Le Chapitre~\ref{chapitre:validation} expose les
protocoles de validation et les résultats expérimentaux. Le
Chapitre~\ref{chapitre:discussion} discute les forces, les limites et les implications
pratiques du système. La conclusion ferme le mémoire et propose des perspectives de
recherche.

\end{introduction}
